\documentclass[11pt,a4paper]{article}
\usepackage[utf8]{inputenc}
\usepackage[T1]{fontenc}
\usepackage[brazilian]{babel}
\usepackage{geometry}
\usepackage{graphicx}
\usepackage{booktabs}
\usepackage{float}
\usepackage{enumitem}
\usepackage{amsmath}
\usepackage{xcolor}
\usepackage{tcolorbox}
\geometry{margin=2.5cm}
\setlength{\parindent}{0pt}
\setlength{\parskip}{0.5em}
\tcbuselibrary{skins,breakable}
\newtcolorbox{nota}[1][]{colback=yellow!10,colframe=orange!60!black,fonttitle=\bfseries,title=#1,breakable}

\begin{document}
\begin{center}
  
\includegraphics[width=0.16\textwidth]{../../../figures_global/unb.png}\\[0.3cm]
  {\large\textbf{GABARITO DO PROFESSOR -- USO RESTRITO}}\\[0.2cm]
  {\large FGA0092 -- Princípios de Comunicação para Engenharia}\\[0.2cm]
  \rule{12cm}{0.8pt}\\[0.2cm]
  {\Large\textbf{Prática 05: Modulação FM}}\\[0.2cm]
  \rule{12cm}{0.8pt}
\end{center}

\section*{Gráfico esperado}
\begin{center}
  \includegraphics[width=0.88\textwidth]{grafico_esperado.png}
\end{center}
\textit{O gráfico acima foi gerado pelo script \texttt{gerar\_grafico.py} desta pasta.}

\section*{Respostas esperadas}

\begin{enumerate}[label=\textbf{Q\arabic*.}, leftmargin=*, itemsep=1em]

  \item \textbf{O que é o desvio de frequência $\Delta f$ e como ele se manifesta no domínio do tempo?}

    \textbf{Resposta esperada:} No FM, $\Delta f$ é o desvio máximo de frequência
    instantânea em relação a $f_c$. Quando $\Delta f$ aumenta, a frequência instantânea
    oscila mais intensamente em torno de $f_c$, e no domínio do tempo a portadora
    aparece com variações mais rápidas de frequência. Para $\Delta f \gg f_m$ (WBFM),
    a forma de onda no tempo assemelha-se a uma senoide de frequência variável;
    para $\Delta f \ll f_m$ (NBFM), visualmente lembra um sinal AM.

    \begin{nota}[Erros comuns]
      Confundir desvio de frequência ($\Delta f$) com variação de amplitude, ou afirmar
      que a amplitude do sinal FM varia com $\Delta f$.
    \end{nota}

  \item \textbf{Aplique a Regra de Carson para calcular a largura de banda para $f_m=1000$\,Hz e $\Delta f=5000$\,Hz. Como a medição experimental se compara?}

    \textbf{Resposta esperada:} Para $f_m=1000$\,Hz e $\Delta f=5000$\,Hz:
    $\beta = \Delta f / f_m = 5$; $B_T = 2(\Delta f + f_m) = 2(5000+1000) = 12000$\,Hz.
    A Regra de Carson é uma aproximação que contém aproximadamente $98\%$ da potência.
    Erros de $10$--$20\%$ são esperados na medição experimental devido à subjetividade
    da medição visual e ao truncamento das componentes de Bessel de ordem alta.

    \begin{nota}[Erros comuns]
      Medir a largura de banda como o intervalo entre os picos mais extremos visíveis
      sem usar critério de potência consistente ($-3$\,dB ou visual padronizado).
    \end{nota}

  \item \textbf{Qual é a diferença espectral entre NBFM e WBFM? Quando cada um é preferível?}

    \textbf{Resposta esperada:} NBFM ($\beta < 1$): espectro com apenas 3 linhas
    principais ($f_c$, $f_c \pm f_m$), similar ao AM; largura de banda $\approx 2f_m$.
    WBFM ($\beta \gg 1$): espectro com muitas linhas em $f_c \pm n f_m$ ponderadas
    por $J_n(\beta)$; largura de banda $\approx 2\beta f_m = 2\Delta f$. WBFM é
    preferível quando a qualidade de áudio é prioritária (radiodifusão FM); NBFM é
    usado em comunicações de voz móvel onde espectro é escasso.

    \begin{nota}[Erros comuns]
      Descrever o WBFM como apenas ``mais potente'' sem mencionar que ele ocupa
      significativamente mais largura de banda que o NBFM.
    \end{nota}

  \item \textbf{Por que o FM oferece maior robustez ao ruído em comparação com o AM? Qual é o ganho teórico para $\beta=5$?}

    \textbf{Resposta esperada:} O FM converte variações de amplitude do ruído em
    variações de frequência que são filtradas pelo discriminador. O ganho de SNR do
    WBFM é $\approx 3\beta^2(\beta+1)$ vezes o do AM com mesma potência. Para
    $\beta=5$: ganho $\approx 3 \times 25 \times 6 = 450 \approx 27$\,dB. O custo é
    $(\beta+1)$ vezes mais espectro que o AM. Abaixo do limiar ($\text{SNR} < 10$\,dB),
    o discriminador colapsa abruptamente.

    \begin{nota}[Erros comuns]
      Afirmar que o FM é ``sempre melhor'' sem mencionar o fenômeno de limiar do FM e
      sem apresentar o cálculo numérico do ganho de SNR.
    \end{nota}

  \item \textbf{Compare experimentalmente NBFM e WBFM em presença de ruído. Os resultados são consistentes com a teoria?}

    \textbf{Resposta esperada:} O WBFM deve suportar níveis de ruído significativamente
    mais altos que o NBFM antes da degradação visível. Se \texttt{noise\_amp\_NBFM}
    $\approx 0{,}1$ e \texttt{noise\_amp\_WBFM} $\approx 0{,}5$ causam degradação
    similar, isso reflete um ganho de SNR de fator 5, consistente com $\beta$ diferindo
    por fator $\sim 5$ entre os casos. O importante é que a tendência seja explicada
    fisicamente com referência à teoria de FM.

    \begin{nota}[Erros comuns]
      Não relacionar os valores de \texttt{noise\_amp} com a previsão teórica de ganho
      de SNR, ou reportar resultados sem comparar com o valor esperado de ganho.
    \end{nota}

\end{enumerate}

\section*{Critérios de correção}
\begin{tabular}{ll}
  \toprule
  Critério & Pontuação \\
  \midrule
  Q1 -- Definição de $\Delta f$ e manifestação no tempo sem confundir com amplitude & 1,0 ponto \\
  Q2 -- Cálculo correto de $\beta$ e $B_T$ pela Regra de Carson                    & 1,0 ponto \\
  Q3 -- Diferença espectral NBFM vs.\ WBFM e contexto de aplicação                & 1,0 ponto \\
  Q4 -- Ganho de SNR com cálculo numérico e menção ao limiar do FM                 & 1,0 ponto \\
  Q5 -- Comparação experimental consistente com tendência teórica                  & 1,0 ponto \\
  \midrule
  \textbf{Total} & \textbf{5,0 pontos} \\
  \bottomrule
\end{tabular}

\bigskip
\begin{nota}[Nota de correção geral]
  Priorizar coerência física e evidência experimental (capturas de tela, tabelas).
  Respostas sem evidência visual devem ter desconto. Não penalizar pequenos erros
  de cálculo se o raciocínio físico estiver correto.
\end{nota}

\end{document}
